\section{Gravity: the Action that Shapes the Mesh}
\label{sec:gravity}

Every other force lives on the mesh. Gravity is the mesh. It is the responsiveness of the geometry of relations to what the relations carry. In Vibe Theory this is not a separate result but a reading of the dynamics already in hand.

The action that drives the causal-set Monte Carlo is the Benincasa-Dowker functional (Benincasa and Dowker, 2010), the discrete analogue of the Einstein-Hilbert action, whose continuum limit is the Ricci scalar curvature integrated over spacetime. It is built from the abundances of small order intervals, and it is the discrete curvature of the note structure.

\begin{principle}[Gravity is the stabilizing action]
The same action whose smeared form makes smooth spacetime a stable phase (Section~\ref{sec:geometry}) is the discrete gravitational action. The dynamics that stabilizes a smooth geometry and the dynamics of gravity are one functional, not two.
\end{principle}

This unifies two things that look separate. In Section~\ref{sec:geometry} the smeared Benincasa-Dowker action was the term that stabilized the manifold phase against the entropic layered background. Here we name what that term is: it is gravity. So in Vibe Theory the stability of a smooth spacetime and the existence of a gravitational action are the same fact, seen once as a phase and once as a law of curvature.

What we have shown is the action and its phase, not the full field equations. We have not derived the Einstein equations as equations of motion, nor a graviton as a propagating excitation. Those are the natural next steps, and they are program rather than result. The standing claim is the one above: the functional that shapes the mesh into a smooth, curved geometry is the discrete Einstein-Hilbert action, and it is the same functional that makes spacetime stable.
